%======================================================================
% Appendix B: Formula sheet
%======================================================================
\chapter{Formula Sheet}\label{app:formulas}

The most important formulas of each chapter, collected for quick reference. Each
heading names the chapter where the formulas are explained, with examples and
conditions of validity. Notation is listed in \cref{app:notation}.

\providecommand{\fsheading}[2]{%
  \par\needspace{4\baselineskip}\medskip
  \noindent{\bfseries\color{MLchapter}#1}\hfill{\footnotesize\cref{#2}}\par
  \nopagebreak\smallskip}

\begin{multicols}{2}
\raggedcolumns
\small
\setlength{\abovedisplayskip}{4pt plus 1pt}
\setlength{\belowdisplayskip}{4pt plus 1pt}
\setlength{\abovedisplayshortskip}{2pt plus 1pt}
\setlength{\belowdisplayshortskip}{2pt plus 1pt}

%---------------------------------------------------------------------
\fsheading{Numbers and algebra}{ch:numbers}
Powers, products and sums:
\begin{align*}
  a^ma^n &= a^{m+n}, \qquad (a^m)^n = a^{mn},\\
  a^{-n} &= \frac{1}{a^n}, \qquad a^{1/n} = \sqrt[n]{a},\\
  (a + b)^2 &= a^2 + 2ab + b^2,\\
  (a + b)(a - b) &= a^2 - b^2,\\
  \sum_{i=1}^n(ax_i + b) &= a\sum_{i=1}^nx_i + nb .
\end{align*}
$\abs{x - a} < r$ means $a - r < x < a + r$. Multiplying an inequality by a
negative number reverses it.

%---------------------------------------------------------------------
\fsheading{Functions}{ch:functions}
Slope of the line through $(x_1, y_1)$ and $(x_2, y_2)$:
$m = \frac{y_2 - y_1}{x_2 - x_1}$.
Roots of $ax^2 + bx + c$ and the vertex:
\[
  x = \frac{-b \pm \sqrt{b^2 - 4ac}}{2a}, \qquad x_{\text{vertex}} = -\frac{b}{2a} .
\]
Exponentials and logarithms:
\begin{align*}
  e^ae^b &= e^{a+b}, \qquad \ln(ab) = \ln a + \ln b,\\
  \ln(a^r) &= r\ln a, \qquad \log_bx = \frac{\ln x}{\ln b} .
\end{align*}
Sigmoid and its inverse, the logit:
\begin{align*}
  \sigma(x) &= \frac{1}{1 + e^{-x}}, \qquad 1 - \sigma(x) = \sigma(-x),\\
  \sigma^{-1}(p) &= \ln\frac{p}{1 - p} .
\end{align*}

%---------------------------------------------------------------------
\fsheading{Geometry and trigonometry}{ch:geometry}
Distance: $d(\vx, \vy) = \sqrt{\sum_i(x_i - y_i)^2}$. Angles: $\pi$ radians $= 180^\circ$.
\begin{align*}
  &\sin^2\theta + \cos^2\theta = 1, \qquad \tan\theta = \frac{\sin\theta}{\cos\theta},\\
  &c^2 = a^2 + b^2 - 2ab\cos C \quad\text{(law of cosines)},\\
  &\sin(\alpha + \beta) = \sin\alpha\cos\beta + \cos\alpha\sin\beta,\\
  &\cos(\alpha + \beta) = \cos\alpha\cos\beta - \sin\alpha\sin\beta .
\end{align*}
Rotation of $(x, y)$ by the angle $\alpha$:
\[
  x' = x\cos\alpha - y\sin\alpha, \qquad y' = x\sin\alpha + y\cos\alpha .
\]

%---------------------------------------------------------------------
\fsheading{Vectors and matrices}{ch:linalg1}
Dot products, norms and matrix algebra:
\begin{align*}
  \vx\cdot\vy &= \textstyle\sum_ix_iy_i = \norm{\vx}\norm{\vy}\cos\theta,\\
  \norm{\vx} &= \sqrt{\vx\cdot\vx}, \qquad \abs{\vx\cdot\vy} \le \norm{\vx}\norm{\vy},\\
  \norm{\vx + \vy} &\le \norm{\vx} + \norm{\vy},\\
  (\mA\mB)_{ij} &= \textstyle\sum_kA_{ik}B_{kj},\\
  (\mA\mB)^\top &= \mB^\top\mA^\top, \qquad (\mA\mB)^{-1} = \mB^{-1}\mA^{-1},\\
  \det(\mA\mB) &= \det\mA\,\det\mB,\\
  \begin{pmatrix} a & b\\ c & d\end{pmatrix}^{-1}
    &= \frac{1}{ad - bc}\begin{pmatrix} d & -b\\ -c & a\end{pmatrix}.
\end{align*}

%---------------------------------------------------------------------
\fsheading{Core linear algebra}{ch:linalg2}
Norms, projections and decompositions:
\begin{align*}
  \norm{\vx}_p &= \Bigl(\sum_i\abs{x_i}^p\Bigr)^{1/p}, \qquad
    \norm{\vx}_\infty = \max_i\abs{x_i},\\
  \proj_{\va}(\vb) &= \frac{\va\cdot\vb}{\va\cdot\va}\,\va,\\
  \mA^\top\mA\hat\vx &= \mA^\top\vb \quad\text{(least squares)},\\
  \mA\vv &= \lambda\vv, \qquad \det(\mA - \lambda\mI) = 0,\\
  \textstyle\sum_i\lambda_i &= \tr\mA, \qquad \textstyle\prod_i\lambda_i = \det\mA,\\
  \mA &= \mQ\mLambda\mQ^\top \quad\text{(symmetric } \mA\text{)},\\
  \mA &= \mU\mSigma\mV^\top \quad\text{(SVD, any } \mA\text{)}.
\end{align*}
$\mA$ is positive definite if $\vx^\top\mA\vx > 0$ for all $\vx \neq \vzero$; for
symmetric $\mA$ this means all eigenvalues are positive.

%---------------------------------------------------------------------
\fsheading{Calculus of one variable}{ch:calculus}
The derivative and its numerical estimate:
\begin{align*}
  f'(x) &= \lim_{h \to 0}\frac{f(x + h) - f(x)}{h}\\
  &\approx \frac{f(x + h) - f(x - h)}{2h} \quad\text{(small } h\text{)}.
\end{align*}
Rules, Taylor approximation and integrals:
\begin{align*}
  (x^n)' &= nx^{n-1}, \qquad (e^x)' = e^x,\\
  (\ln x)' &= \frac1x, \qquad (\sin x)' = \cos x,\\
  (\cos x)' &= -\sin x,\\
  (fg)' &= f'g + fg',\\
  \Bigl(\frac fg\Bigr)' &= \frac{f'g - fg'}{g^2},\\
  \bigl(f(g(x))\bigr)' &= f'(g(x))\,g'(x),\\
  \sigma' &= \sigma(1 - \sigma),\\
  f(a + h) &\approx f(a) + f'(a)\,h + \tfrac12f''(a)\,h^2,\\
  \int_a^bf'(x)\,\dd x &= f(b) - f(a),\\
  \int_{-\infty}^{\infty}e^{-x^2}\,\dd x &= \sqrt\pi .
\end{align*}

%---------------------------------------------------------------------
\fsheading{Multivariable calculus}{ch:multivariable}
Gradients, chain rule and approximation:
\begin{align*}
  \nabla f &= \Bigl(\pderiv{f}{x_1}, \dots, \pderiv{f}{x_n}\Bigr),\\
  D_{\vu}f &= \nabla f\cdot\vu,\\
  \deriv{}{t}f(\vx(t)) &= \textstyle\sum_i\pderiv{f}{x_i}\deriv{x_i}{t},\\
  \mat J_{\vf\circ\vg} &= \mat J_{\vf}\,\mat J_{\vg},\\
  f(\vx + \vh) &\approx f(\vx) + \nabla f(\vx)^\top\vh + \tfrac12\vh^\top\mH\vh,\\
  \nabla(\va^\top\vx) &= \va,\\
  \nabla(\vx^\top\mA\vx) &= (\mA + \mA^\top)\vx,\\
  \nabla\norm{\mA\vx - \vb}^2 &= 2\mA^\top(\mA\vx - \vb).
\end{align*}
At a critical point ($\nabla f = \vzero$): positive definite Hessian means a
minimum, negative definite a maximum, mixed signs a saddle.

%---------------------------------------------------------------------
\fsheading{Optimisation}{ch:optimization}
Gradient descent: $\vtheta \leftarrow \vtheta - \eta\nabla L(\vtheta)$.\\
SGD: $\vtheta \leftarrow \vtheta - \eta\frac1b\sum_{i\in B}\nabla\ell_i(\vtheta)$.\\
Momentum: $\vv \leftarrow \beta\vv + \vg$, then $\vtheta \leftarrow \vtheta - \eta\vv$.\\
Newton: $\vtheta \leftarrow \vtheta - \mH^{-1}\nabla L(\vtheta)$.\par
Adam, with gradient $\vg_t$ at step $t$ (componentwise):
\begin{align*}
  \vm_t &= \beta_1\vm_{t-1} + (1 - \beta_1)\vg_t,\\
  \vs_t &= \beta_2\vs_{t-1} + (1 - \beta_2)\vg_t^2,\\
  \hat\vm_t &= \frac{\vm_t}{1 - \beta_1^t}, \qquad \hat\vs_t = \frac{\vs_t}{1 - \beta_2^t},\\
  \vtheta_t &= \vtheta_{t-1} - \eta\,\frac{\hat\vm_t}{\sqrt{\hat\vs_t} + \epsilon} .
\end{align*}
Ridge regression: $(\mX^\top\mX + \lambda\mI)\vw = \mX^\top\vy$.\\
Lagrange multipliers: $\nabla f = \lambda\nabla g$ and $g = 0$.\\
Convexity: for all $\lambda \in [0, 1]$,
\[
  f(\lambda\vx + (1 - \lambda)\vy) \le \lambda f(\vx) + (1 - \lambda)f(\vy).
\]

%---------------------------------------------------------------------
\fsheading{Probability}{ch:probability}
Rules of probability, expectation and variance:
\begin{align*}
  P(A^c) &= 1 - P(A),\\
  P(A \cup B) &= P(A) + P(B) - P(A \cap B),\\
  P(A \mid B) &= \frac{P(A \cap B)}{P(B)},\\
  P(A \mid B) &= \frac{P(B \mid A)\,P(A)}{P(B)} \quad\text{(Bayes)},\\
  P(A) &= \textstyle\sum_iP(A \mid B_i)\,P(B_i),\\
  \E[X] &= \textstyle\sum_xx\,p(x),\\
  \E[aX + b] &= a\E[X] + b,\\
  \Var(X) &= \E[X^2] - (\E[X])^2,\\
  \Var(aX + b) &= a^2\Var(X),\\
  \Cov(X, Y) &= \E[XY] - \E[X]\E[Y],\\
  \rho &= \frac{\Cov(X, Y)}{\sigma_X\sigma_Y},\\
  \Var(X + Y) &= \Var X + \Var Y + 2\Cov(X, Y).
\end{align*}
Binomial: $P(X = k) = \binom nk p^k(1 - p)^{n-k}$, mean $np$, variance $np(1 - p)$.
Normal: $p(x) = \frac{1}{\sigma\sqrt{2\pi}}e^{-(x - \mu)^2/(2\sigma^2)}$; $68\%$, $95\%$
and $99.7\%$ lie within $1$, $2$ and $3$ standard deviations of the mean.
Multivariate normal:
\[
  p(\vx) = \frac{\exp\bigl(-\frac12(\vx - \vmu)^\top\mSigma^{-1}(\vx - \vmu)\bigr)}
  {\sqrt{(2\pi)^d\det\mSigma}} .
\]

%---------------------------------------------------------------------
\fsheading{Statistics}{ch:statistics}
Estimation and inference:
\begin{align*}
  \bar x &= \tfrac1n\textstyle\sum_ix_i, \qquad
    s^2 = \tfrac{1}{n-1}\textstyle\sum_i(x_i - \bar x)^2,\\
  \operatorname{SE}(\bar X) &= \sigma/\sqrt n, \qquad
    \text{95\% CI: } \bar x \pm 1.96\,s/\sqrt n,\\
  \MSE(\hat\theta) &= \Bias(\hat\theta)^2 + \Var(\hat\theta),\\
  \hat\theta_{\mathrm{MLE}} &= \argmax_\theta\textstyle\sum_i\ln p(x_i \mid \theta),\\
  \hat\theta_{\mathrm{MAP}} &= \argmin_\theta\bigl(-\ln p(\mathcal{D} \mid \theta) - \ln p(\theta)\bigr).
\end{align*}
Coin: $\hat\theta = h/n$. Normal: $\hat\mu = \bar x$, $\hat\sigma^2 = \frac1n\sum_i(x_i - \bar x)^2$.
Gaussian noise gives squared error; Bernoulli labels give binary cross-entropy
$-\sum_i[y_i\ln q_i + (1 - y_i)\ln(1 - q_i)]$.
Expected test error $=$ bias$^2$ $+$ variance $+$ noise.

%---------------------------------------------------------------------
\fsheading{Deep learning}{ch:dl}
Layers, softmax and information theory:
\begin{align*}
  \vz^{(l)} &= \mW^{(l)}\va^{(l-1)} + \vb^{(l)},\\
  \va^{(l)} &= \varphi(\vz^{(l)}),\\
  \softmax(\vz)_k &= \frac{e^{z_k}}{\sum_je^{z_j}},\\
  L = -\ln q_c &\;\Longrightarrow\; \pderiv{L}{\vz} = \vq - \vy,\\
  H(\vp) &= -\textstyle\sum_kp_k\log p_k,\\
  H(\vp, \vq) &= -\textstyle\sum_kp_k\log q_k,\\
  D_{\KL}(\vp\,\Vert\,\vq) &= \textstyle\sum_kp_k\log\frac{p_k}{q_k}\\
  &= H(\vp, \vq) - H(\vp) \ge 0 .
\end{align*}
For $\vz = \mW\vx + \vb$ with $\vdelta = \pderiv{L}{\vz}$:
$\pderiv{L}{\mW} = \vdelta\vx^\top$, $\pderiv{L}{\vx} = \mW^\top\vdelta$,
$\pderiv{L}{\vb} = \vdelta$. Backpropagation:
\[
  \vdelta^{(l)} = \bigl(\mW^{(l+1)\top}\vdelta^{(l+1)}\bigr) \odot \varphi'(\vz^{(l)}).
\]
Attention: $\softmax\bigl(\mQ\mK^\top/\sqrt d\bigr)\mV$.
\end{multicols}
